\scanpage{4}
2. \uline{Généralités sur certains foncteurs de modules}.

Le présent numéro pourrait passer également au \deleted{\unsure{nº|Nº}} Chap $\unsure{0|O}_{\mathrm{III}}$, qui contiendrait aussi certains des sorites du par.7 .

\uline{Proposition} 2.1. Soient $\underline C$ une catégorie abélienne, $\underline G$ une partie de $\mathrm{Ob}\,\underline C$, telle que tout objet de $\underline C$ admette une suite de composition finie dont les \deleted{quotients}\inserted{facteurs} sont isomorphes à des quotients d'objets $M\handwritten{\in}\underline G$. \deleted{\unsure{Al}} Soit $\underline C'$ une seconde catégorie abélienne.

(i) Soit $T:\underline C\handwritten{\to}\underline{C'}$ un foncteur \deleted{add} exact à droite, pour que $T$ soit nul, il f et s qu'il soit nul sur les $M\handwritten{\in}\underline G$.

(ii) Soit $(T^i)$ un foncteur cohomologique de $\underline C$ dans $\underline C'$, tel que $T^i=0$ pour $i\handwritten{>}i_{\unsure{o|0}}$, et soit $n$ un entier. Pour qu'on ait $T^i=0$ pour $i\handwritten{>}n$, il f. et s. que pour tout $M\handwritten{\in}\underline G$, on ait $T^i(M)=0$ pour $i\handwritten{>}n$.

Dém. (i) est trivial, pour (ii) on procède par récurrence \inserted{descendante} sur $n$, c'est trivial pour $n\handwritten{\geq}i_o$, et pour $n$ quelconque, supposant que c'est fait pour $n+1$, on note que si les $T^i(M)$ ($M\handwritten{\in}\underline G$, $i\handwritten{>1+}n$) sont nuls, alors par l'hy de réc on a $T^i=0$ pour $i\handwritten{>}n+1$, donc $T^{n\handwritten{+1}}$ est exact à droite, et d'après (i) il est nul, cqfd.

\uline{Corollaire} \deleted{\unsure{3}}2.2. Sous les conditions de 2.1. soit $(T^i)$ un foncteur cohomologique contravariant de $\underline C$ dans $\underline C'$, tel que $T^i=0$ pour $i\handwritten{<}i_o$, et soit $n$ un entier. Pour qu'on ait $T^i=0$ pour $i\handwritten{<}\deleted{\unsure{m}}\handwritten{n}$, il f et s que pour tout $M\handwritten{\in}\underline G$, on ait $T^i(M)=0$ pour $i\handwritten{<}n$.

Il suffit d'appliquer 2.1. (ii) à $\underline C$ et ${\underline C'}^{o}$. Pour qu'il n'y ait pas de jaloux, il faudrait aussi répéter (i). Dans le présent paragraphe, ce sera sous cette forme que no\deleted{\unsure{j}}us utiliserons 2.1. (mais il aurait été canularesque d'énoncer 2.1. avec des foncteurs contravariants).

\uline{Corollaire} 2.3. Soient $\underline C$, $\underline C'$ des catégories a\deleted{\unsure{v}}béliennes, $\underline G$ une partie de $\mathrm{Ob}\,\underline C$, on suppose que les objets de $\underline C$ sont noethériens (i.e. satisfont la condition des chaines\handwritten{\unsure{/}}ascendantes pour les sous-objets), et que $\underline G$ est un système de générateurs de $\underline C$ (i.e. que pour tout $\deleted{\unsure{A|E}}F\handwritten{\ne}0$ dans $\underline C$, il existe un $M\handwritten{\in}\underline C$ et un homomorphisme non nul $M\handwritten{\to}F$). Alors les conclusions de 2.1. (i)(ii) et 2.2. (i)(ii) sont valables.

Il suffit de vérifier l\deleted{es}\handwritten{'} hypothèses de 2.1.\inserted{\handwritten{ sur $\underline G$}}, pour ceci soit $F\handwritten{\in}\mathrm{Ob}\,\underline C$, \deleted{\unsure{si F}} soit $F'$ un sous-objet maximal de $F$ admettant une suite de composition finie comme dite dans 2.1. (un tel $F'$ existe, car $0\handwritten{\subset}F$ satisfait à l'hypothèse de 2.1., et on applique l'hyp noeth sur $F$). Je dis que $F'=F$, en \deleted{\unsure{effe}}
% archive sheet number: 3

